\scanpage{176}
Je dis qu'on a alors une \deleted{équivalence}\inserted{``\unsure{homéomorphie}''} (heuristique !)
\[
(292)\qquad\check N^{d'_0}_{d_*,d''_0}\mid D^*_{d_*}\simeq\unsure{\mathcal V}^{d'_0}_{d_*,d''_0}
\]
qui induit, grâce à l'hom
\[
(293)\qquad\check N^{d'_0}_{d_*,d''_0}\mid D^*_{d_*}\longrightarrow\check N^*_{d_*;(d'_0,d''_0)}\mid D^*_{d_*}
\]
qui \uline{est une équivalence d'homotopie}

une équivalence d'homotopie
\[
(294)\qquad\boxed{\begin{gathered}
\unsure{\mathcal V}^{d'_0}_{d_*,d''_0}\simeq\check N^*_{d_*;d'_0,d''_0}\mid D^*_{d_*}\\[1em]
\text{équivalence d'homotopie}
\end{gathered}}
\]
qui a le seul inconvénient de ne pas correspondre à une flèche \uline{canonique} dans un sens ni dans l'autre, donc (pour le moment) d'avoir une valeur purement heuristique. Retenons cependant la
\margin{175}
